\documentclass[10pt,a4paper]{article}
\usepackage[margin=24mm]{geometry}
\usepackage[T1]{fontenc}
\usepackage{lmodern,microtype,booktabs,amsmath,graphicx}
\usepackage[numbers,sort&compress]{natbib}
\usepackage{url}
\usepackage[hidelinks,pdfauthor={},pdftitle={Online appendix}]{hyperref}
\setlength{\parindent}{0pt}
\setlength{\parskip}{4pt}
\title{Online appendix\\[3pt]\large Benchmark corpora for detecting LLM-era bots}
\author{}
\date{}
\begin{document}
\maketitle
\vspace{-2em}
This appendix holds what the paper points to. Every number in it is reproduced by the code in this
repository (\texttt{scripts/run\_all.sh}) and checked against \texttt{results/expected/}.

\section{How the ten corpora were found}\label{app:search}
The search ran on 20 August 2026 with web search and the OpenAlex API, after exploratory
  searches earlier that month.
It followed three axes.
First, keyword variants for how the field names these bots (``LLM-driven'', ``LLM-powered'',
  ``AI-powered'' and ``generative agent'', each with ``social bot'', ``detection'', ``benchmark''
  or ``dataset'').
Second, searches by dataset name, using each corpus already found as an entry point.
Third, a forward citation chase through OpenAlex from two anchors, Fox8-23 (57 citing works at
  the time of the search) and BotSim-24 (12).
The chase found the most; the keyword axis found the least, because the field's naming is
  unstable.
Each axis stopped when a query returned nothing new.
Which corpora a paper evaluates on was read from its full text, since abstract pages do not name
  them.

We kept every account-labelled corpus of social-media bots the search returned that either carries a
  published account-level detection result from a model that uses content or was released to
  study a campaign or information operation.
Corpora that label texts rather than accounts, such as MultiSocial, are out of scope
  (Sect.~2.1 of the paper).
We did not keep the exact wording of every keyword query or per-query screening counts, so this
  appendix records the procedure and its stopping rule, not a screening count.

\section{The ten corpora}\label{app:corpora}
That search returns the ten corpora in Table~1 of the paper.
Six have published account-level detection results from models that use content.
The other four were built mainly to study campaigns and information operations rather than
  detection.

\paragraph{Most of the ten lack one of the three.}
\emph{Three hold no LLM-written posts.}
TwiBot-20 \citep{feng2021twibot20} contains accounts crawled outwards from seed users in politics,
  business, entertainment and sports.
TwiBot-22 \citep{feng2022twibot22} uses a larger crawl of the same kind; both release the posts.
Twitter IO \citep{seckin2024labeled} contains accounts Twitter removed as state-linked influence
  operations.
All three were collected before ChatGPT's release in November 2022.

\emph{Three have no human class.}
CMV \citep{jaidka2026howfar} releases 33 undisclosed LLM accounts that argued in a Reddit debate
  forum, all of them the experiment's bots \citep{cmvmods2025disclosure}.
Yard Signs \citep{linvill2024yardsigns} contains only accounts from a 2024 political campaign
  on X that replied to ordinary users, with no human comparison class.
Masquerade-23 \citep{li2025masquerade} is text from Chirper.ai, a platform whose users are all
  bots, so on its own it has no human class; we use it below in rebuilt form.

\emph{Two release no posts.}
MGTAB \citep{shi2023mgtab} contains expert-annotated accounts from a Twitter debate and releases
  extracted features rather than posts.
Bluesky \citep{meng2026navigating} contains 10,610 accounts from the platform Bluesky, with a label and
  source for each account but no posts.

For Table~1 of the paper, MGTAB, TwiBot-20 and TwiBot-22 each used one crawl and labelled
  the accounts afterwards.
Bluesky is partly LLM-written: 460 of its 4,215 labelled bots are marked as generative.
MGTAB releases embeddings rather than posts; Bluesky releases account labels and their sources.

\section{Baselines}\label{app:baselines}
On Fox8-23 and BotSim-24 our runs use the accounts with at least 50 usable posts (20 characters or more), and score every account out of fold under five-fold cross-validation.
We use cross-validation because neither corpus ships a split: the published results use a random
  held-out split that each paper draws in its own code.
On Masquerade-rebuilt the published scores are for a three-way task that adds TwiBot-20's
  pre-LLM bots, so they are not directly comparable with our two-class result.

\paragraph{The simple baseline.}
Where the tables say ``bag of words'' the model is TF-IDF over the account's own posts with
  a minimum document frequency of two, followed by logistic regression, scored out of fold
  under five-fold cross-validation.
The whole pipeline is fitted inside each training fold, and no two accounts share a document.
The Fox8-23 and BotSim-24 bot-versus-human runs join each account's 50 messages;
  Sect.~
ef{app:io} gives the separate Twitter IO settings.
The human-source comparisons on Fox8-23 join 20 messages, so these two analyses do not use the
  same amount of text.

Masquerade-rebuilt is scored on the text contrast it creates, Chirper.ai posts against TwiBot-20's
  human accounts, with 5,158 accounts per class.

We report the bag-of-words result as AUC everywhere except Sect.~2.2 of the paper, where we
  also compute F1 to compare with published F1 scores.
The papers do not state their F1 thresholds; the released TRACE-Bot and BotDMM code takes the
  more probable class, a cut at 0.5, which we also use.
In the 50-message condition, any threshold between the two classes gives the same result:
  the out-of-fold scores do not overlap,
  the lowest bot scoring 0.737 against the highest human 0.471 on Fox8-23 and 0.704 against 0.240
  on BotSim-24, so the AUC is exactly 1 and every threshold between them gives the same F1.
The bag of words also needs little text: 20 randomly selected posts concatenated and then
  truncated to 25 words in total still give F1 0.866 on Fox8-23 and 0.897 on BotSim-24.
This truncated condition does overlap, and its F1 does depend on the cut.
% NEW: how a printed number was computed
TwiBot-22's bag of words is scored on the 6,920 labelled accounts with at least one English reply to an English
  post that is not in a thread of their own; each account's replies are joined and cut to 20,000 characters.

\paragraph{The attribute baseline.}
A single account attribute uses the same five-fold procedure on one feature, so its direction is
  learned on the training folds rather than chosen with the labels.
The attribute itself is picked by whichever scores best, and the shuffled-label null repeats
  that choice, which is why those nulls sit slightly above 0.5.
For each pair of Fox8-23's human sources, we choose the best of seven attributes, such as the
  follower count or the length of the bio.
On BotSim-24 the best attribute is whether the account has a bio, a non-empty profile
  description: all 1,000 generated accounts have one, against 278 of the 1,907 human accounts
  (14.6\%).

Twitter IO is scored one campaign at a time, each with whichever attribute separates its own
  accounts best.
Of its 27 campaigns, 25 have at least 20 accounts on each side and can be scored.
Creation date gives the highest AUC in 21 of those 25; Table~1 of the paper gives the
  range of the best attribute's AUC across the scored campaigns.
% NEW: how a printed number was computed
Above 10,000 accounts (TwiBot-20, MGTAB, TwiBot-22) the quantile transform behind the attribute baseline is fitted
  on a random subsample of 10,000 accounts; we fix its seed.

\paragraph{The Fox8-23 bag of words without the self-revealing phrases.}
The phrases are ``as an AI language model'' and seven variants, such as ``I'm an AI'',
  ``OpenAI'' and ``ChatGPT''; no human account in the cohort uses any of them.
On the same 500 bots and 500 humans, deleting the phrases from every message leaves AUC and F1 at
  1.0000.
Dropping the accounts whose 50 scored messages contain one keeps 222 bots, and the scores stay at
  1.0000.
Dropping those whose full released history contains one keeps 31 bots; AUC stays at 1.0000, and
  against 31 humans drawn at random, 100 times, F1 averages 0.9989.
Shuffled labels give 0.47 to 0.51 in the first three runs.

\paragraph{Masquerade-rebuilt built as the detector papers describe it.}
Table~1 of the paper scores Masquerade-rebuilt as the text contrast it creates.
Its published scores are for a different task: BotDMM, LLM-BotGuard and BotEvo classify humans,
  TwiBot-20's traditional bots and LLM-driven bots.
To compare on that task, we built the corpus as BotDMM and LLM-BotGuard describe it, with our own
  draws where they are silent, and scored the same bag of words (Table~\ref{tab:masquerade}).
Against humans, the LLM-driven bots separate at F1 0.9900, as on Fox8-23 and BotSim-24.
On the three-way task the same model reaches macro F1 0.8194, below the published 0.89 to 0.94.
What it separates less well is TwiBot-20's traditional bots from its humans (all bots against
  humans, AUC 0.8976), the same difficulty as Table~1's TwiBot-20 row.
Table~1's attribute cell reads ``accounts are TwiBot-20's'' for the same reason:
  the LLM-driven accounts are TwiBot-20 bot accounts with replaced text.
Under the BotDMM recipe and on Table~1's fields, the followers/friends ratio
  separates them from humans at
  0.7495, against 0.7485 for all of TwiBot-20's bots, and no account count separates them from the
  traditional bots (0.48 to 0.51).

\begin{table}[t]\centering\small
\setlength{\tabcolsep}{4pt}
\resizebox{\linewidth}{!}{\begin{tabular}{llcccccc}
\toprule
construction & task and split & accuracy & macro F1 & weighted F1 & MCC & AUC & F1 \\
\midrule
BotDMM recipe & three-way, 70/20/10 test & 0.8115 & 0.8194 & 0.8089 & 0.7084 & & \\
 & three-way, five-fold CV & 0.8150 & 0.8227 & 0.8117 & 0.7141 & & \\
 & LLM-driven bots vs humans, 70/20/10 test & 0.9806 & & & & 0.9991 & 0.9730 \\
 & LLM-driven bots vs humans, five-fold CV & 0.9927 & & & & 0.9995 & 0.9900 \\
 & all bots vs humans, 70/20/10 test & 0.8014 & & & & 0.8939 & 0.8151 \\
 & all bots vs humans, five-fold CV & 0.8082 & & & & 0.8976 & 0.8219 \\
\addlinespace
LLM-BotGuard recipe & three-way, 70/20/10 test & 0.8191 & 0.8024 & 0.8108 & 0.7224 & & \\
 & three-way, five-fold CV & 0.8316 & 0.8179 & 0.8251 & 0.7410 & & \\
 & LLM-driven bots vs humans, 70/20/10 test & 0.9886 & & & & 0.9994 & 0.9857 \\
 & LLM-driven bots vs humans, five-fold CV & 0.9926 & & & & 0.9996 & 0.9907 \\
 & all bots vs humans, 70/20/10 test & 0.8132 & & & & 0.9068 & 0.8228 \\
 & all bots vs humans, five-fold CV & 0.8238 & & & & 0.9103 & 0.8342 \\
\midrule
published, BotDMM & three-way & 0.9542 & 0.9362 & & 0.9025 & & \\
published, LLM-BotGuard & three-way & 0.8850 & 0.8851 & & 0.8212 & & \\
published, LLM-BotGuard & bots vs humans & 0.8817 & & & & & \\
published, BotEvo & three-way & & 0.8930 & & & & \\
\bottomrule
\end{tabular}}
\caption{Built by the detector papers' recipe, Masquerade-rebuilt still lets a bag of words
  separate LLM-driven bots from humans, but the bag of words falls below the published three-way scores.
BotDMM replaces the tweets and profile description of 3,000 TwiBot-20 bot accounts with those of a
  Masquerade-23 (Chirper.ai) account and relabels them LLM-driven; LLM-BotGuard replaces the tweets of 3,500;
  BotEvo follows BotDMM.
The papers do not state which accounts were drawn, how much of a timeline is replaced, or which F1
  average is printed, so our rows are our own draws under their stated recipe, not their files.
Our rows use TF-IDF over an account's words and logistic regression, fitted on the 70\% split and
  scored on the 10\% test split, or scored out of fold under five-fold cross-validation; the last
  column is the bot class's F1 at a 0.5 cut.
The published rows are the papers' full models on their own draws.}
\label{tab:masquerade}
\end{table}

\paragraph{A bag of words fitted on one corpus and scored on the other.}
The transfer results use the text cohorts of Sect.~\ref{app:baselines}, 500 bots and 500 humans for
  Fox8-23 and 329 and 408 for BotSim-24, 50 messages each joined into one document.
The bag of words is fitted on every account of one corpus and scored on every account of the other,
  so no test account is seen in training and no cross-validation is needed.
Fitted on Fox8-23 it scores 0.998 on BotSim-24, and fitted on BotSim-24 it scores 0.890 on Fox8-23;
  resampling the scored accounts 1,000 times puts these between 0.996 and 1.000 and between 0.869
  and 0.909.
Deleting ``as an AI language model'' from both corpora changes neither.
Document length alone gives 0.781 and 0.744, and training on shuffled labels 0.51 and 0.49.
Both sets of bots were written by OpenAI models, which may explain part of this transfer.
In both corpora the humans were also gathered apart from the bots, so these transfer scores do not
  separate shared writing features from collection differences.

\section{Collection routes}\label{app:routes}
\begin{table}[t]\centering\small
\setlength{\tabcolsep}{4pt}
\resizebox{\linewidth}{!}{\begin{tabular}{lrrcrrrr}
\toprule
                        &          &        & single-class & route-only & majority &      &     \\
corpus, and recorded route   & accounts & routes & routes     & accuracy   & rate & gain & $V$ \\
\midrule
Fox8-23, a source dataset     &   2,280 &  5 & 5 of 5  & 1.0000 & 0.5000 & $+0.5000$ & 1.0000 \\
BotSim-24, generated or scraped &  2,907 &  2 & 2 of 2  & 1.0000 & 0.6560 & $+0.3440$ & 1.0000 \\
Masquerade-rebuilt, Chirper.ai or TwiBot-20 & 10,316 &  2 & 2 of 2  & 1.0000 & 0.5000 & $+0.5000$ & 1.0000 \\
Twitter IO, the archive or the control search & 319,686 &  2 & 2 of 2  & 1.0000 & 0.9509 & $+0.0491$ & 1.0000 \\
Bluesky, a recorded source    & 10,610 &  9 & 8 of 9  & 0.9955 & 0.6027 & $+0.3927$ & 0.9907 \\
TwiBot-20, a seed domain      &  11,826 & 15 & 0 of 15 & 0.6553 & 0.5572 & $+0.0982$ & 0.3712 \\
\bottomrule
\end{tabular}}
\caption{The collection route alone perfectly predicts the label in all three LLM-era candidates.
A route is where an account comes from and what brings it into the corpus, such as a search, an
  earlier study or a community report.
The table covers every corpus whose account routes are recorded or fixed by construction.
Route-only accuracy is the accuracy of predicting the most common class within each route;
  majority rate is that of always predicting the largest class; gain is their difference.
Cram\'er's $V$ measures route--label association (0: independent; 1: route determines label).
Only in TwiBot-20 does every route contain both classes.
In BotSim-24, Masquerade-rebuilt and Twitter IO the route is the class by construction, so their
  accuracy of 1 is true by definition, not measured.}
\label{tab:route}
\end{table}

In Table~\ref{tab:route}, BotSim-24, Masquerade-rebuilt and Twitter IO each use a separate route
  for each class by construction.
Bluesky records what brought each account in, such as a purchase or a community report; eight
  of its nine sources contain only one class.
\citet{meng2026navigating} report excluding features used in constructing their Bluesky dataset
  from their classifiers.
\textbf{We measure how well the recorded collection route alone predicts the label}, Bluesky included
  (Table~\ref{tab:route}), and we require one route for both classes (Sect.~2.3 of the paper).

\section{Every Twitter IO campaign, scored against its own controls}\label{app:io}
Twitter IO is the one family that releases its own controls, so each campaign is scored without
  us choosing the human side.
Simple baselines still separate the classes: the median campaign has bag-of-words AUC 0.9665
  (Table~\ref{tab:io}).
Creation date is the best single attribute in 21 of the 25 scored campaigns, with median AUC 0.91.
The creation-date result has two possible explanations, and the release does not separate them.
Campaign accounts may really be registered in bursts, which is behaviour a detector should
  use, or the difference may come from how the controls were gathered, since they were chosen by
  hashtag and day and not by creation date.

\paragraph{How the two classes were gathered.}
The operation's accounts come from Twitter's archive with their whole timelines.
The controls are accounts that used one of the operation's hashtags on the same day, with at most
  100 of their posts from those days, and the authors note that a control's posts cover a short
  period while an operation account's can span years \citep{seckin2024labeled}.
Across all 27 campaigns, the operation's accounts have more posts than their controls in 26,
  at a median ratio of 7.0.
Their posts span a median of 137 days, against 0 for controls, and the most common language
  differs between the classes in 12 campaigns.

\paragraph{Scoring.}
Table~\ref{tab:io} scores the 25 campaigns with enough accounts in both classes.
The account-level baseline uses up to 200 posts per account, drops reposts, keeps messages of
  at least 20 characters without a language filter, caps each class at 4,000 accounts, and uses
  five-fold cross-validation.
Bangladesh (11 accounts) and Russia\_3 (5) fall below the floor of 20 accounts per class and
  are not scored.

\paragraph{Matching the two classes exactly by hashtag and day.}
The rule both classes share is a post under a given hashtag on a given day, so the match is made
  per message (Table~\ref{tab:iomatch}).
All four rows use posts in the operation's most common language, at most five posts per account,
  and the same number of posts per class.
The first row draws posts at random.
The second keeps only posts under a hashtag and day that both classes used.
The third gives both classes exactly the same number of posts under each hashtag and day, and the
  fourth also deletes hashtags, mentions and URLs from the text and drops any text that occurs more
  than once.
Folds are grouped by account, the AUC is computed inside each fold and averaged, and each cell is
  the mean of ten draws.
Four of the 27 campaigns keep at least 100 exactly matched posts from at least 20 accounts in each
  class in all ten draws, and only those are scored.
One Ecuador account appears in both classes; all of its posts are excluded.
Shuffling the labels by account gives 0.49 to 0.51 in every cell.

Matching frequencies matters.
An account-level check using three posts per account matched which hashtags and days appeared,
  but not how often.
AUC fell in 2 of 10 campaigns and rose at the median, from 0.873 to 0.911.
The exact message-level match controls those frequencies in both classes.

\begin{table}[t]\centering\small
\begin{tabular}{lcccr}
\toprule
campaign & bag of words & length & null & accounts, campaign / control \\
\midrule
Qatar           & 0.8087 & 0.7127 & 0.5065 & 28 / 4,000 \\
China\_2        & 0.8698 & 0.8273 & 0.5214 & 179 / 4,000 \\
Russia\_4       & 0.9033 & 0.7987 & 0.5041 & 24 / 4,000 \\
Venezuela\_2    & 0.9058 & 0.9737 & 0.5187 & 33 / 3,002 \\
Russia\_2       & 0.9143 & 0.6644 & 0.4814 & 351 / 1,892 \\
Thailand        & 0.9177 & 0.5363 & 0.5285 & 301 / 1,386 \\
Iran\_2         & 0.9491 & 0.6437 & 0.4915 & 289 / 4,000 \\
Iran\_3         & 0.9535 & 0.7137 & 0.5224 & 135 / 1,659 \\
China\_1        & 0.9548 & 0.7130 & 0.4948 & 683 / 4,000 \\
Russia\_5       & 0.9556 & 0.7578 & 0.4683 & 48 / 4,000 \\
Cuba\_2         & 0.9605 & 0.8255 & 0.5144 & 422 / 528 \\
Ghana/Nigeria   & 0.9607 & 0.8685 & 0.4352 & 57 / 885 \\
Ecuador         & 0.9665 & 0.6862 & 0.4965 & 768 / 4,000 \\
Iran\_1         & 0.9666 & 0.6985 & 0.4951 & 644 / 4,000 \\
Spain           & 0.9672 & 0.8250 & 0.4761 & 210 / 1,067 \\
Iran\_6         & 0.9673 & 0.8774 & 0.4849 & 206 / 4,000 \\
Catalonia       & 0.9686 & 0.6565 & 0.5458 & 47 / 1,723 \\
UAE             & 0.9758 & 0.8196 & 0.4961 & 3,432 / 4,000 \\
Egypt/UAE       & 0.9759 & 0.8461 & 0.4709 & 236 / 285 \\
Russia\_1       & 0.9763 & 0.7297 & 0.5088 & 3,123 / 4,000 \\
Cuba\_1         & 0.9800 & 0.8852 & 0.4871 & 377 / 4,000 \\
Iran\_4         & 0.9821 & 0.6075 & 0.4899 & 2,468 / 4,000 \\
Armenia         & 0.9837 & 0.7253 & 0.5121 & 31 / 1,384 \\
Venezuela\_1    & 0.9928 & 0.9136 & 0.5005 & 577 / 4,000 \\
Iran\_5         & 0.9938 & 0.3054 & 0.5505 & 104 / 889 \\
\bottomrule
\end{tabular}
\caption{A bag of words exceeds AUC 0.90 in 23 of 25 campaigns with no LLM-written posts.
Each Twitter IO campaign is compared with its released controls at account level, using five-fold
  cross-validation (median AUC 0.9665; range 0.8087--0.9938).
``Length'' uses only the character count of joined posts; ``null'' uses shuffled labels.
Below-chance AUCs retain their direction.}
\label{tab:io}
\end{table}

\begin{table}[t]\centering\small
\begin{tabular}{lcccc}
\toprule
posts compared & Cuba\_1 & Ecuador & Russia\_1 & Spain \\
\midrule
drawn at random                  & 0.710 & 0.819 & 0.783 & 0.867 \\
hashtag and day used by both     & 0.655 & 0.844 & 0.815 & 0.872 \\
same number per hashtag and day  & 0.471 & 0.673 & 0.569 & 0.683 \\
\quad and no tags, no repeated text & 0.495 & 0.632 & 0.581 & 0.612 \\
\addlinespace
posts per class                  & 158 & 307 & 195 & 137 \\
\bottomrule
\end{tabular}
\caption{Recovering a shared route lowers message-level AUC in all four Twitter IO campaigns.
Median bag-of-words AUC falls from 0.801 to 0.621 when bots and controls contribute equal numbers
  of posts per shared hashtag and day.
Rows progress from random posts to shared hashtag--day pairs, then equal class counts within
  each pair; the final row also removes hashtags, mentions, URLs and repeated text.
AUCs average ten draws with cross-validation folds grouped by account.}
\label{tab:iomatch}
\end{table}
% NEW: how the matched rows were sampled
In the last two rows of Table~\ref{tab:iomatch} every hashtag and day holds the same number of bot and control
  posts in every draw.

\section{Collecting a corpus, in full}\label{app:collect}
The paper's Sect.~5 gives the short form of this advice; this is the long form it is shortened from.

\paragraph{Where to get it.}
Most corpora come from a platform's API or from scraping.
Both see only what is public and still online: the platform's own records stay out of reach, and
  removed accounts are gone.
The EU's Digital Services Act gives researchers another way to obtain data
  \citep{eu2022dsa}.
Under its Article 40, a researcher can ask a very large online platform for data to study
  systemic risks.
Bot activity is one of them: Article 34(2) names ``inauthentic use or automated exploitation of
  the service''.
The request can cover public data (Article 40(12)) and, for vetted researchers, non-public data
  (Article 40(4), with access procedures in force since October 2025
  \citep{eu2025dataaccess}).
In either case a researcher can ask for bots and humans collected by one route, and for the
  records needed to check it.
What a platform hands over is settled case by case.

\paragraph{What to keep.}
We recommend keeping everything the platform's terms and the study's ethics approval allow.
A post not kept now may be gone later, and collecting only what is still online introduces
  survivorship bias.
Three things are easy to lose and hard to reconstruct:
\begin{itemize}
\item \textbf{Whole account histories.} As many posts per account as possible, for both classes,
  let an evaluation compare a few posts with a longer history. An analysis can subsample later,
  but it cannot recover posts that were never kept.
\item \textbf{Each account's route.} A record of how every account entered the corpus, including
  any search terms used to find it, lets researchers check whether the route predicts the label
  and test a detector on accounts from a route held out of training (Sect.~2.3 of the paper).
\item \textbf{The post each reply answers.} Its identifier and text let replies to the same posts
  be collected by one rule for both classes, which gives reply bots a shared route
  (Sect.~2.3 of the paper). It also lets future detectors use the conversation, not the reply alone.
\end{itemize}
\bibliographystyle{plainnat}
\bibliography{references}
\end{document}
